% SOURCE of the appendix (single-column submission only; scripts/build_paper.py fills it into appendix.tex and
% main_single.tex inputs it after the references). Same rules as body.template.tex: every number is a double-brace
% placeholder read from a result file, and every analysis added in the revision is labelled (exploratory).
% The body refers to these sections as appendix A1 to A7; main_single.tex numbers them A1, A2, ...

\section{Training settings}\label{app:A1}
All rollout-trained models of the F18 and five-split pools were trained with the settings in Table~\ref{tab:app-train}.
The optimiser is {{hp_optimiser}}; gradient clipping: {{hp_clip}}. The loss has {{hp_n_loss_terms}} non-zero terms,
with weights {{hp_loss_weights}}. These settings were checked to be identical in the training records of all
{{hp_n_records}} parent models of the five-split pool. The recurrent network has two GRU layers of {{hp_hidden}} units
and reads {{hp_M}} past steps.

\begin{table}[h]
\centering
\caption{What differs between the rollout-trained parent, the teacher-forced models and the $N{=}1$ models. A blank means
the setting equals the parent's. Records compared: {{hp_tf_n_records}} teacher-forced and {{hp_n1_n_records}} $N{=}1$ models.}
\label{tab:app-train}
\footnotesize
\fittable{\begin{tabular}{@{}lccc@{}}
\toprule
Setting & Rollout-trained parent & Teacher-forced & $N{=}1$ \\
\midrule
Learning rate / weight decay / batch & {{hp_lr}} / {{hp_wd}} / {{hp_batch}} & & \\
Iterations & {{hp_iters}} & & \\
Loss target & own predictions fed back & recorded states fed in & own predictions \\
Steps predicted per window ($N$) & {{hp_parent_N}} & {{hp_parent_N}} & {{hp_n1_N}} \\
\bottomrule
\end{tabular}}
\end{table}

Other F18 models differ from the parent in one setting each. The longer-trained models ran {{hp_10k_iters}}
iterations (the teacher-forced ones as well). The ensembles have {{hp_ens_n}} members. The capacity-matched control has
{{hp_w124_hidden}} hidden units. The likelihood model replaces the squared error by a Gaussian likelihood, and two
further models train on contaminated or duplicated data. The compressed copies are made from the finished parents without
further training: 16-bit and 8-bit weights, removal of the smallest weights over the whole network (10, 20, 30 and
50\%), and removal of units so that the hidden width is 192 or 128.

\section{What was known before the registration}\label{app:A2}
Section 2 of the registration (\texttt{PREREG.md}, supplementary material) states what we had read before any model was run through this study's code.
For the growth rate $\lambda_1$ on the F18 pool, an earlier study found Spearman {{pre_ng_R_rho}} {{pre_ng_R_ci}}
with the 368-step error over {{pre_ng_R_n}} rollout-trained models, and {{pre_ng_RT_rho}} {{pre_ng_RT_ci}} over
{{pre_ng_RT_n}} rollout-trained and teacher-forced ones. On the five splits, $\lambda_1$ ranked {{pre_icmai_all_n}} models at
{{pre_icmai_all_rho}} {{pre_icmai_all_ci}}, and the {{pre_icmai_roll_n}} rollout-trained ones among them at
{{pre_icmai_roll_rho}} {{pre_icmai_roll_ci}} (not settled). The mean kinematic residual ranked {{pre_icp_n}} models at
{{pre_icp_rho}} {{pre_icp_ci}}, much of it through the difference between training methods. Earlier studies also
showed that teacher forcing beats rollout training at one step but loses at 368 steps, and that moderate pruning
without fine-tuning loses skill at long horizons.
Open were the one-step and 8-step errors computed from every position of the held-out data, the correlation of any
short-horizon error with any long-horizon error, and the growth rate or residual of any compressed copy. The raw
arrays of the earlier growth-rate study were used for one consistency check only, and no statistic is taken from them.

\section{Horizon-matched and split-half checks}\label{app:A3}
Table~\ref{tab:app-horizon} gives the Spearman correlation with the 100-step error $E_{100}$ when the growth rate
and the kinematic residual are computed over a window of 8, 24, 100 or 368 steps (exploratory). The one-step and 8-step
errors are shown for comparison. Over 8 steps the growth rate's correlation is {{a1_F18_lam1_8_rho}} on F18 and
{{a1_X5_lam1_8_rho}} on X5, against {{a1_F18_lam1_100_rho}} and {{a1_X5_lam1_100_rho}} over 100 steps.
With both checks computed over 8 steps, the residual's lead over the 8-step error was {{a1_F18_rj8_e8_delta}}
{{a1_F18_rj8_e8_ci}} on F18 and {{a1_X5_rj8_e8_delta}} {{a1_X5_rj8_e8_ci}} on X5; over 100 steps it was
{{a1_F18_rj100_e8_delta}} {{a1_F18_rj100_e8_ci}} on F18 and {{a1_X5_rj100_e8_delta}} {{a1_X5_rj100_e8_ci}} on X5 (all exploratory).

\begin{table}[h]
\centering
\caption{Spearman correlation of each check with $E_{100}$ by the window the check is computed over, for F18 (4 held-out
trajectories; all simulators, trained, compressed) and X5 (20 trajectories) (exploratory).}
\label{tab:app-horizon}
\footnotesize
\setlength{\tabcolsep}{4pt}
\fittable{\begin{tabular}{@{}lcccc@{}}
\toprule
Check (window) & F18, all & F18, trained & F18, compressed & X5 \\
\midrule
One-step error $E_1$ & {{a1t_F18all_E1}} & {{a1t_F18tr_E1}} & {{a1t_F18cp_E1}} & {{a1t_X5_E1}} \\
8-step error $E_8$ & {{a1t_F18all_E8}} & {{a1t_F18tr_E8}} & {{a1t_F18cp_E8}} & {{a1t_X5_E8}} \\
\midrule
$\lambda_1$ (8 steps) & {{a1t_F18all_lam1_8}} & {{a1t_F18tr_lam1_8}} & {{a1t_F18cp_lam1_8}} & {{a1t_X5_lam1_8}} \\
$\lambda_1$ (24 steps) & {{a1t_F18all_lam1_24}} & {{a1t_F18tr_lam1_24}} & {{a1t_F18cp_lam1_24}} & {{a1t_X5_lam1_24}} \\
$\lambda_1$ (100 steps) & {{a1t_F18all_lam1_100}} & {{a1t_F18tr_lam1_100}} & {{a1t_F18cp_lam1_100}} & {{a1t_X5_lam1_100}} \\
$\lambda_1$ (368 steps) & {{a1t_F18all_lam1_368}} & {{a1t_F18tr_lam1_368}} & {{a1t_F18cp_lam1_368}} & {{a1t_X5_lam1_368}} \\
\midrule
$\bar r_J$ (8 steps) & {{a1t_F18all_rJ_8}} & {{a1t_F18tr_rJ_8}} & {{a1t_F18cp_rJ_8}} & {{a1t_X5_rJ_8}} \\
$\bar r_J$ (24 steps) & {{a1t_F18all_rJ_24}} & {{a1t_F18tr_rJ_24}} & {{a1t_F18cp_rJ_24}} & {{a1t_X5_rJ_24}} \\
$\bar r_J$ (100 steps) & {{a1t_F18all_rJ_100}} & {{a1t_F18tr_rJ_100}} & {{a1t_F18cp_rJ_100}} & {{a1t_X5_rJ_100}} \\
$\bar r_J$ (368 steps) & {{a1t_F18all_rJ_368}} & {{a1t_F18tr_rJ_368}} & {{a1t_F18cp_rJ_368}} & {{a1t_X5_rJ_368}} \\
\bottomrule
\end{tabular}}
\end{table}

Table~\ref{tab:app-split} computes each check on one half of the trajectories and the target on the other half
(exploratory), beside the correlation when both use the same half. The largest drop is {{a3_max_drop}}.

\begin{table}[h]
\centering
\caption{Split-half check (exploratory): Spearman correlation with $E_{100}$ when the check and the target come from
different halves of the held-out trajectories, with its 95\% interval, and when both come from the same half.}
\label{tab:app-split}
\footnotesize
\setlength{\tabcolsep}{4pt}
\fittable{\begin{tabular}{@{}lcccc@{}}
\toprule
 & F18, other half & F18, same half & X5, other half & X5, same half \\
\midrule
$E_1$ & {{a3_F18_E1_rho}} {{a3_F18_E1_ci}} & {{a3_F18_E1_same}} & {{a3_X5_E1_rho}} {{a3_X5_E1_ci}} & {{a3_X5_E1_same}} \\
$E_8$ & {{a3_F18_E8_rho}} {{a3_F18_E8_ci}} & {{a3_F18_E8_same}} & {{a3_X5_E8_rho}} {{a3_X5_E8_ci}} & {{a3_X5_E8_same}} \\
$\lambda_1$ & {{a3_F18_lam1_rho}} {{a3_F18_lam1_ci}} & {{a3_F18_lam1_same}} & {{a3_X5_lam1_rho}} {{a3_X5_lam1_ci}} & {{a3_X5_lam1_same}} \\
$\bar r_J$ & {{a3_F18_rJ_rho}} {{a3_F18_rJ_ci}} & {{a3_F18_rJ_same}} & {{a3_X5_rJ_rho}} {{a3_X5_rJ_ci}} & {{a3_X5_rJ_same}} \\
\bottomrule
\end{tabular}}
\end{table}

\section{Wider intervals, the five-split pool without split \{1,8\}, and Holm's correction}\label{app:A4}
All of this section is exploratory. The registered intervals on correlations resample models (within each split on
X5). The wider intervals resample parent models and held-out trajectories together, in two levels, with {{a2_ndraws}} draws. They are {{a2_median_ratio}} times as wide as the
registered ones (median; {{a2_median_ratio_t2}} for the correlations of Table~\ref{tab:checks} and {{a2_median_ratio_tests}} for
the registered differences; range {{a2_ratio_min}} to {{a2_ratio_max}}). Of the {{a2_n_tests}} registered verdicts,
{{a2_n_hold}} hold and {{a2_n_lost}} do not.

\begin{table}[h]
\centering
\caption{Table~\ref{tab:checks}'s correlations with $E_{100}$, with the registered interval (where one was computed) and
the two-level 95\% interval (exploratory); for X5 also without the split \{1, 8\}, which repeats F18's held-out pair (exploratory).}
\label{tab:app-two}
\footnotesize
\setlength{\tabcolsep}{3pt}
\fittable{\begin{tabular}{@{}llcccc@{}}
\toprule
Pool & Check & Registered & Registered 95\% & Two-level 95\% & Without \{1, 8\} \\
\midrule
F18, all & $E_1$ & {{q1_rho_e1}} & {{rho_E1_E100_ci}} & {{a2_t2_f18_all_e1_ci}} & \\
 & $E_8$ & {{q2_rho_e8}} & {{rho_E8_E100_ci}} & {{a2_t2_f18_all_e8_ci}} & \\
 & $\lambda_1$ & {{q1_rho_lam}} & {{rho_lam1_E100_ci}} & {{a2_t2_f18_all_lam1_ci}} & \\
 & $\bar r_J$ & {{q4_1_rho_new}} & {{rho_rJ_E100_ci}} & {{a2_t2_f18_all_rj_ci}} & \\
F18, trained & $E_1$ & {{x2_tc_E1}} & {{x2_tc_E1_ci}} & {{a2_t2_f18_trained_e1_ci}} & \\
 & $E_8$ & {{x2_tc_E8}} & {{x2_tc_E8_ci}} & {{a2_t2_f18_trained_e8_ci}} & \\
 & $\lambda_1$ & {{x2_tc_lam1}} & {{x2_tc_lam1_ci}} & {{a2_t2_f18_trained_lam1_ci}} & \\
 & $\bar r_J$ & {{x2_tc_rJ}} & {{x2_tc_rJ_ci}} & {{a2_t2_f18_trained_rj_ci}} & \\
F18, compressed & $E_1$ & {{x2_c_E1}} & & {{a2_t2_f18_compressed_e1_ci}} & \\
 & $E_8$ & {{x2_c_E8}} & & {{a2_t2_f18_compressed_e8_ci}} & \\
 & $\lambda_1$ & {{x2_c_lam1}} & & {{a2_t2_f18_compressed_lam1_ci}} & \\
 & $\bar r_J$ & {{x2_c_rJ}} & & {{a2_t2_f18_compressed_rj_ci}} & \\
X5 & $E_1$ & {{q3_1_rho_base}} & & {{a2_t2_x5_e1_ci}} & {{a2b_X5_E1_rho}} {{a2b_X5_E1_ci}} \\
 & $E_8$ & {{q3_8_rho_base}} & & {{a2_t2_x5_e8_ci}} & {{a2b_X5_E8_rho}} {{a2b_X5_E8_ci}} \\
 & $\lambda_1$ & {{q3_1_rho_new}} & & {{a2_t2_x5_lam1_ci}} & {{a2b_X5_lam1_rho}} {{a2b_X5_lam1_ci}} \\
 & $\bar r_J$ & {{x1_rJ}} (E) & & {{a2_t2_x5_rj_ci}} & {{a2b_X5_rJ_rho}} {{a2b_X5_rJ_ci}} \\
\bottomrule
\end{tabular}}
\end{table}

Without split \{1, 8\}, $\lambda_1$ minus $E_1$ is {{a2b_q3_vs_e1_delta}} {{a2b_q3_vs_e1_ci}}, $\lambda_1$ minus $E_8$ is
{{a2b_q3_vs_e8_delta}} {{a2b_q3_vs_e8_ci}}, and $\bar r_J$ minus $E_8$ is {{a2b_rj_vs_e8_delta}} {{a2b_rj_vs_e8_ci}}.

Table~\ref{tab:app-tests} lists the registered differences. Holm's correction keeps the chance of any false finding among
the {{a2c_n_family}} secondary tests of Q3--Q6 at {{a2c_alpha}}, and {{a2c_n_survive}} of them survive. The readings of Q1 and
Q2 at 24 and 368 steps, study FT's secondary tests and every exploratory analysis are outside this family and
uncorrected. Its $p$-values come from a
Monte Carlo resampling, so a value close to {{a2c_alpha}} is itself uncertain (``near'' marks the tests within two
standard errors of the threshold).

\begin{table}[h]
\centering
\caption{Registered differences in Spearman correlation (first minus second check; Q1 and Q2 at 97.5\%, the rest at 95\%),
with the registered and the two-level interval, whether the registered verdict still holds, and, for the seven secondary tests of
Q3--Q6, Holm's adjusted $p$ with whether it survives and whether it is near the threshold. Two-level and Holm columns are exploratory.}
\label{tab:app-tests}
\footnotesize
\setlength{\tabcolsep}{3pt}
\fittable{\begin{tabular}{@{}lccccc@{}}
\toprule
Test & Estimate & Registered & Two-level & Holds & Holm $p$ (survives; near) \\
\midrule
Q1: $\lambda_1-E_1$ at 100 steps, F18 & {{q1_delta}} & {{q1_ci975}} & {{a2_q1_ci}} & {{a2_q1_holds}} & -- \\
Q2: $\lambda_1-E_8$ at 100 steps, F18 & {{q2_delta}} & {{q2_ci975}} & {{a2_q2_ci}} & {{a2_q2_holds}} & -- \\
Q1 at 24 steps & {{q1_e24_delta}} & {{q1_e24_ci}} & {{a2_q1_e24_ci}} & {{a2_q1_e24_holds}} & -- \\
Q1 at 368 steps & {{q1_e368_delta}} & {{q1_e368_ci}} & {{a2_q1_e368_ci}} & {{a2_q1_e368_holds}} & -- \\
Q2 at 24 steps & {{q2_e24_delta}} & {{q2_e24_ci}} & {{a2_q2_e24_ci}} & {{a2_q2_e24_holds}} & -- \\
Q2 at 368 steps & {{q2_e368_delta}} & {{q2_e368_ci}} & {{a2_q2_e368_ci}} & {{a2_q2_e368_holds}} & -- \\
Q3: $\lambda_1-E_1$, five splits & {{q3_1_delta}} & {{q3_1_ci}} & {{a2_q3_vs_e1_ci}} & {{a2_q3_vs_e1_holds}} & {{a2c_q3_vs_e1_pholm}} ({{a2c_q3_vs_e1_survives}}; {{a2c_q3_vs_e1_near}}) \\
Q3: $\lambda_1-E_8$, five splits & {{q3_8_delta}} & {{q3_8_ci}} & {{a2_q3_vs_e8_ci}} & {{a2_q3_vs_e8_holds}} & {{a2c_q3_vs_e8_pholm}} ({{a2c_q3_vs_e8_survives}}; {{a2c_q3_vs_e8_near}}) \\
Q4: $\bar r_J-E_1$, F18 & {{q4_1_delta}} & {{q4_1_ci}} & {{a2_q4_vs_e1_ci}} & {{a2_q4_vs_e1_holds}} & {{a2c_q4_vs_e1_pholm}} ({{a2c_q4_vs_e1_survives}}; {{a2c_q4_vs_e1_near}}) \\
Q4: $\bar r_J-E_8$, F18 & {{q4_8_delta}} & {{q4_8_ci}} & {{a2_q4_vs_e8_ci}} & {{a2_q4_vs_e8_holds}} & {{a2c_q4_vs_e8_pholm}} ({{a2c_q4_vs_e8_survives}}; {{a2c_q4_vs_e8_near}}) \\
Q5: $\lambda_1-E_1$, no one-step-trained & {{q5_1_delta}} & {{q5_1_ci}} & {{a2_q5_vs_e1_ci}} & {{a2_q5_vs_e1_holds}} & {{a2c_q5_vs_e1_pholm}} ({{a2c_q5_vs_e1_survives}}; {{a2c_q5_vs_e1_near}}) \\
Q5: $\lambda_1-E_8$, no one-step-trained & {{q5_8_delta}} & {{q5_8_ci}} & {{a2_q5_vs_e8_ci}} & {{a2_q5_vs_e8_holds}} & {{a2c_q5_vs_e8_pholm}} ({{a2c_q5_vs_e8_survives}}; {{a2c_q5_vs_e8_near}}) \\
Q6: growth-rate shift of pruned copies & {{q6_mean}} & {{q6_ci}} & {{a2_q6_ci}} & {{a2_q6_holds}} & {{a2c_q6_pholm}} ({{a2c_q6_survives}}; {{a2c_q6_near}}) \\
\bottomrule
\end{tabular}}
\end{table}

\section{Compression and the growth rate}\label{app:A5}
Table~\ref{tab:app-comp} gives, for every compression variant of the {{q6_nparents}} five-split parents, the shift in
growth rate, how many copies have $\lambda_1>0$, and the skill retained $K$ (the share of the parent's advantage over
predicting no change that the copy keeps). In all, {{n_unstable_removed}} of the {{n_removed_copies}} copies with
weights or units removed have $\lambda_1>0$, all of them among the {{n_wide_copies}} copies pruned by 50\% or reduced in
width.

\begin{table}[h]
\centering
\caption{Compressed copies of the five-split parents: mean growth rate $\lambda_1$ of the copies, mean shift from the
parent (95\% interval), copies with $\lambda_1>0$, and skill retained $K$ at 8, 100 and 368 steps (negative: worse than predicting
no change).}
\label{tab:app-comp}
\footnotesize
\setlength{\tabcolsep}{3pt}
\fittable{\begin{tabular}{@{}lccccccc@{}}
\toprule
Variant & $\lambda_1$ (1/s) & Shift & Copies $\lambda_1>0$ & $K_8$ & $K_{100}$ & $K_{368}$ \\
\midrule
16-bit weights & {{lam_fp16}} & {{lamshift_fp16}} {{lamshift_fp16_ci}} & {{lampos_fp16}} & {{K_fp16_h8}} & {{K_fp16_h100}} & {{K_fp16_h368}} \\
8-bit weights & {{lam_int8w}} & {{lamshift_int8w}} {{lamshift_int8w_ci}} & {{lampos_int8w}} & {{K_int8w_h8}} & {{K_int8w_h100}} & {{K_int8w_h368}} \\
10\% pruned & {{lam_u10}} & {{lamshift_u10}} {{lamshift_u10_ci}} & {{lampos_u10}} & {{K_u10_h8}} & {{K_u10_h100}} & {{K_u10_h368}} \\
20\% pruned & {{lam_u20}} & {{lamshift_u20}} {{lamshift_u20_ci}} & {{lampos_u20}} & {{K_u20_h8}} & {{K_u20_h100}} & {{K_u20_h368}} \\
30\% pruned & {{lam_u30}} & {{lamshift_u30}} {{lamshift_u30_ci}} & {{lampos_u30}} & {{K_u30_h8}} & {{K_u30_h100}} & {{K_u30_h368}} \\
50\% pruned & {{lam_u50}} & {{lamshift_u50}} {{lamshift_u50_ci}} & {{lampos_u50}} & {{K_u50_h8}} & {{K_u50_h100}} & {{K_u50_h368}} \\
Width 192 & {{lam_s192}} & {{lamshift_s192}} {{lamshift_s192_ci}} & {{lampos_s192}} & {{K_s192_h8}} & {{K_s192_h100}} & {{K_s192_h368}} \\
Width 128 & {{lam_s128}} & {{lamshift_s128}} {{lamshift_s128_ci}} & {{lampos_s128}} & {{K_s128_h8}} & {{K_s128_h100}} & {{K_s128_h368}} \\
\bottomrule
\end{tabular}}
\end{table}

An earlier registered study of the same recipe, with 500 fine-tuning iterations, is summarised in
Table~\ref{tab:app-perc}. Fine-tuned narrow copies are compared with the equally fine-tuned full-width model, as in the
body; the control row fine-tunes the unreduced model in the same way. Its last column compares fine-tuning
with training from scratch at the same width. In a registered test of PyTorch's 8-bit dynamic quantisation (weights
and activations; the release's relative error), the model kept {{perc_int8_K}} {{perc_int8_K_ci}} of its 368-step
skill.

\begin{table}[h]
\centering
\caption{Skill retained $K$ of fine-tuned narrow copies at 100 and 368 steps in the first run, its replication and
the pooled runs (95\% intervals); the unreduced control (pooled); and the registered label and difference for fine-tuning
against training from scratch at 368 steps.}
\label{tab:app-perc}
\footnotesize
\setlength{\tabcolsep}{3pt}
\fittable{\begin{tabular}{@{}llccc@{}}
\toprule
Reading & Width & $K_{100}$ & $K_{368}$ & Fine-tuned vs scratch (368 steps) \\
\midrule
First run & 192 & {{perc_K_first_w192_h100}} {{perc_K_first_w192_h100_ci}} & {{perc_K_first_w192_h368}} {{perc_K_first_w192_h368_ci}} & {{perc_Q4_first_w192}} {{perc_Q4_first_w192_ci}} \\
 & 128 & {{perc_K_first_w128_h100}} {{perc_K_first_w128_h100_ci}} & {{perc_K_first_w128_h368}} {{perc_K_first_w128_h368_ci}} & {{perc_Q4_first_w128}} {{perc_Q4_first_w128_ci}} \\
Replication & 192 & {{perc_K_rep_w192_h100}} {{perc_K_rep_w192_h100_ci}} & {{perc_K_rep_w192_h368}} {{perc_K_rep_w192_h368_ci}} & {{perc_Q4_rep_w192}} {{perc_Q4_rep_w192_ci}} \\
 & 128 & {{perc_K_rep_w128_h100}} {{perc_K_rep_w128_h100_ci}} & {{perc_K_rep_w128_h368}} {{perc_K_rep_w128_h368_ci}} & {{perc_Q4_rep_w128}} {{perc_Q4_rep_w128_ci}} \\
Pooled & 192 & {{perc_K_pool_w192_h100}} {{perc_K_pool_w192_h100_ci}} & {{perc_K_pool_w192_h368}} {{perc_K_pool_w192_h368_ci}} & {{perc_Q4_pool_w192}} {{perc_Q4_pool_w192_ci}} \\
 & 128 & {{perc_K_pool_w128_h100}} {{perc_K_pool_w128_h100_ci}} & {{perc_K_pool_w128_h368}} {{perc_K_pool_w128_h368_ci}} & {{perc_Q4_pool_w128}} {{perc_Q4_pool_w128_ci}} \\
Control (unreduced) & & {{perc_ctrl_pool_h100}} {{perc_ctrl_pool_h100_ci}} & {{perc_ctrl_pool_h368}} {{perc_ctrl_pool_h368_ci}} & -- \\
\bottomrule
\end{tabular}}
\end{table}

\section{Run-time signals}\label{app:A6}
A check that ranks whole simulators need not flag individual bad steps. Earlier analyses of these models, on
the same held-out trajectories of the five splits, found that none of the label-free signals is a usable alarm for a
step that is worse than holding the last state. On the rollout-trained ensembles, a combined kinematic-consistency
score reached an area under the ROC curve (AUROC; 0.5 is chance) of {{icp_auc_phys}} {{icp_auc_phys_ci}}, ensemble
disagreement {{icp_auc_dis}} and a plain step counter {{icp_auc_counter}}. Averaged over models of all three kinds, the growth
rate read step by step (over the last eight steps) reached {{ng_q3_auc}}, below the step counter's {{ng_q3_counter}}
(difference {{ng_q3_delta}} {{ng_q3_ci}}; registered: \textsc{worse than counter}). Re-synchronising the simulator with the recorded state when the
kinematic score fired gave {{icp_h7_excess}} {{icp_h7_ci}} more 368-step error than the same number of
re-synchronisations on a fixed schedule (registered: \textsc{schedule better}). Physics helps as a repair only where it is exact: rescaling the predicted
gravity direction to unit length changed the 368-step error by {{icp_G}} {{icp_G_ci}}, while forcing the joint
angles onto the approximate integration rule multiplied it by {{icp_J_x}}.

\section{Study FT: fine-tuning the pruned copies}\label{app:A7}
% A7: the study registered in prereg/PREREG_FT.md before any fine-tuned model existed (R3): design, arms, decision
% rules and results from scripts/ft_stats.py (results/ft_stats.json); no departure was recorded in ANOMALIES.md.
\textbf{Design.} Study FT asked whether a short fine-tuning, with the pruned weights held at zero, restores the
long-horizon skill that magnitude pruning removes; its rules and analysis code were committed before any fine-tuned pruned model existed. The parents were the {{ft_nparents}} rollout-trained five-split
models with seeds {{ft_seeds}}, the seeds for which an earlier study had fine-tuned unpruned controls. Each parent's
smallest 30\% of weight-matrix entries ({{ft_pruned_u30}}) were set to zero, and the copy was fine-tuned for
{{ft_iters}} iterations by the controls' protocol (the parent's training episodes, rollout loss, settings and seed,
fresh optimiser state), with the pruned entries reset to zero after every step; 50\% pruning ({{ft_pruned_u50}}
entries) was a secondary arm. The parent, the copies before and after fine-tuning and the control were all scored
on this machine by the measures of Sec.~\ref{sec:setup}, on each parent's own split, giving {{ft_ntraj}}
trajectories with seeds averaged within each. Two primary questions at 100 steps, with 97.5\% intervals from
{{ft_ndraws}} trajectory resamples, were registered: FT1, the change in $K$ against the original parent from
fine-tuning (interval above 0: fine-tuning restores skill); and FT2, $K$ of the fine-tuned copy against the control,
with the earlier study's thresholds (lower end above 90\%: nearly as skilled; above 50\%: keeps most; upper end below
50\%: loses most).

\textbf{Time rule and checks.} The 50\% arm was to run only if a probe, run before any fine-tuning, projected the
whole study at no more than {{ft_cap_h}}~h; it projected {{ft_probe_total_h}}~h, so both arms ran. The
{{ft_nruns}} runs took {{ft_wall_total_h}}~h in all ({{ft_wall_min_min}} to {{ft_wall_max_min}} minutes each, mean
{{ft_wall_mean_min}}), and scoring the {{ft_n_scored}} models took {{ft_score_min}} minutes. All four registered checks
passed for every run ({{ft_cft1_pass}} each): the pruned entries were exactly zero after training and after
reloading; every loss trace was finite; the settings equalled the parent's training record; and this machine's
100-step errors for the parent and the unretrained copy matched the stored values within {{ft_cft4_maxrel}}
(relative; the registered tolerance was $10^{-3}$). No arm was excluded, and no bootstrap draw was undefined.

\textbf{Results.} Both primary questions were settled (Table~\ref{tab:app-ft}): fine-tuning restored the 100-step
skill (FT1), and the fine-tuned copies were nearly as skilled as the controls (FT2). The secondary readings at 24 and
368 steps, and at 50\% pruning, agree. Against the original parent, $K$ above 100\% mostly reflects the fine-tuning
itself: the control, fine-tuned without pruning, reaches {{ft_K_CTRL_vs_P_h100}} {{ft_K_CTRL_vs_P_h100_ci}} at 100
steps (descriptive). The unretrained term of FT1 uses seeds {{ft_seeds}} only, hence {{ft_FT1_Knoft}} rather than
Table~\ref{tab:decisions}'s {{K_u30_h100}} (seeds 0--4). After fine-tuning, the one-step reading of $K$ could not be
separated from the 100-step one (secondary; not settled). Unretrained pruning raised the kinematic residual and did
not measurably move the growth rate; after fine-tuning, both were slightly below the parent's (secondary). Table~\ref{tab:app-ft-desc} gives $K$ of every arm (descriptive). As everywhere in this paper, this is
one robot, one trotting gait, simulated data and recorded commands, with no policy trained.

\begin{table}[h]
\centering
\caption{Study FT: registered questions. Primary: 97.5\% intervals (a family of two). Secondary: 95\% intervals,
not corrected for multiplicity. $d$: difference in $K$ (percentage points). $\Delta$: mean change from the parent
($\lambda_1$ in 1/s; $\bar r_J$ in units of the recording's own residual). U30, U50: pruned 30\% or 50\%, not
retrained; F: then fine-tuned.}
\label{tab:app-ft}
\footnotesize
\setlength{\tabcolsep}{3pt}
\fittable{\begin{tabular}{@{}lllcl@{}}
\toprule
 & Question & Quantity & Estimate [interval] & Verdict \\
\midrule
Primary & FT1, 100 steps & $d=K(\mathrm{U30F})-K(\mathrm{U30})$; {{ft_FT1_Kft}} against {{ft_FT1_Knoft}} & {{ft_FT1_d}} {{ft_FT1_ci}} & \textsc{{{ft_FT1_verdict}}} \\
 & FT2, 100 steps & $K$ of U30F against the control & {{ft_FT2_K}} {{ft_FT2_ci}} & \textsc{{{ft_FT2_verdict}}} \\
\midrule
Secondary & FT1, 24 steps & $d$; {{ft_FT1_h24_Kft}} against {{ft_FT1_h24_Knoft}} & {{ft_FT1_h24_d}} {{ft_FT1_h24_ci}} & \textsc{{{ft_FT1_h24_verdict}}} \\
 & FT2, 24 steps & $K$ of U30F against the control & {{ft_FT2_h24_K}} {{ft_FT2_h24_ci}} & \textsc{{{ft_FT2_h24_verdict}}} \\
 & FT1, 368 steps & $d$; {{ft_FT1_h368_Kft}} against {{ft_FT1_h368_Knoft}} & {{ft_FT1_h368_d}} {{ft_FT1_h368_ci}} & \textsc{{{ft_FT1_h368_verdict}}} \\
 & FT2, 368 steps & $K$ of U30F against the control & {{ft_FT2_h368_K}} {{ft_FT2_h368_ci}} & \textsc{{{ft_FT2_h368_verdict}}} \\
 & FT1, 50\%, 100 steps & $d$; {{ft_FT1_u50_h100_Kft}} against {{ft_FT1_u50_h100_Knoft}} & {{ft_FT1_u50_h100_d}} {{ft_FT1_u50_h100_ci}} & \textsc{{{ft_FT1_u50_h100_verdict}}} \\
 & FT2, 50\%, 100 steps & $K$ of U50F against the control & {{ft_FT2_u50_h100_K}} {{ft_FT2_u50_h100_ci}} & \textsc{{{ft_FT2_u50_h100_verdict}}} \\
 & One-step reading & $K(\mathrm{U30F})$ at 1 step minus at 100 steps & {{ft_onestep_U30F_d}} {{ft_onestep_U30F_ci}} & \textsc{{{ft_onestep_U30F_verdict}}} \\
 & Growth rate & $\Delta\lambda_1$, U30F (parent {{ft_change_lam1_U30F_P_P}}) & {{ft_change_lam1_U30F_P_mean}} {{ft_change_lam1_U30F_P_ci}} & \textsc{{{ft_change_lam1_U30F_P_verdict}}} \\
 &  & $\Delta\lambda_1$, U30 & {{ft_change_lam1_U30_P_mean}} {{ft_change_lam1_U30_P_ci}} & \textsc{{{ft_change_lam1_U30_P_verdict}}} \\
 & Kinematic residual & $\Delta\bar r_J$, U30F (parent {{ft_change_rJ_U30F_P_P}}) & {{ft_change_rJ_U30F_P_mean}} {{ft_change_rJ_U30F_P_ci}} & \textsc{{{ft_change_rJ_U30F_P_verdict}}} \\
 &  & $\Delta\bar r_J$, U30 & {{ft_change_rJ_U30_P_mean}} {{ft_change_rJ_U30_P_ci}} & \textsc{{{ft_change_rJ_U30_P_verdict}}} \\
\bottomrule
\end{tabular}}
\end{table}

\begin{table}[h]
\centering
\caption{Study FT, descriptive: skill retained $K$ of each arm against the original parent and against the control
(95\% intervals, {{ft_ntraj}} trajectories, seeds {{ft_seeds}}). CTRL: the parent fine-tuned the same way without
pruning.}
\label{tab:app-ft-desc}
\footnotesize
\setlength{\tabcolsep}{3pt}
\fittable{\begin{tabular}{@{}llcccc@{}}
\toprule
Arm & Against & 1 step & 24 steps & 100 steps & 368 steps \\
\midrule
U30 & parent & {{ft_K_U30_vs_P_h1}} {{ft_K_U30_vs_P_h1_ci}} & {{ft_K_U30_vs_P_h24}} {{ft_K_U30_vs_P_h24_ci}} & {{ft_K_U30_vs_P_h100}} {{ft_K_U30_vs_P_h100_ci}} & {{ft_K_U30_vs_P_h368}} {{ft_K_U30_vs_P_h368_ci}} \\
U30F & parent & {{ft_K_U30F_vs_P_h1}} {{ft_K_U30F_vs_P_h1_ci}} & {{ft_K_U30F_vs_P_h24}} {{ft_K_U30F_vs_P_h24_ci}} & {{ft_K_U30F_vs_P_h100}} {{ft_K_U30F_vs_P_h100_ci}} & {{ft_K_U30F_vs_P_h368}} {{ft_K_U30F_vs_P_h368_ci}} \\
U50 & parent & {{ft_K_U50_vs_P_h1}} {{ft_K_U50_vs_P_h1_ci}} & {{ft_K_U50_vs_P_h24}} {{ft_K_U50_vs_P_h24_ci}} & {{ft_K_U50_vs_P_h100}} {{ft_K_U50_vs_P_h100_ci}} & {{ft_K_U50_vs_P_h368}} {{ft_K_U50_vs_P_h368_ci}} \\
U50F & parent & {{ft_K_U50F_vs_P_h1}} {{ft_K_U50F_vs_P_h1_ci}} & {{ft_K_U50F_vs_P_h24}} {{ft_K_U50F_vs_P_h24_ci}} & {{ft_K_U50F_vs_P_h100}} {{ft_K_U50F_vs_P_h100_ci}} & {{ft_K_U50F_vs_P_h368}} {{ft_K_U50F_vs_P_h368_ci}} \\
CTRL & parent & {{ft_K_CTRL_vs_P_h1}} {{ft_K_CTRL_vs_P_h1_ci}} & {{ft_K_CTRL_vs_P_h24}} {{ft_K_CTRL_vs_P_h24_ci}} & {{ft_K_CTRL_vs_P_h100}} {{ft_K_CTRL_vs_P_h100_ci}} & {{ft_K_CTRL_vs_P_h368}} {{ft_K_CTRL_vs_P_h368_ci}} \\
\midrule
U30F & control & {{ft_K_U30F_vs_CTRL_h1}} {{ft_K_U30F_vs_CTRL_h1_ci}} & {{ft_K_U30F_vs_CTRL_h24}} {{ft_K_U30F_vs_CTRL_h24_ci}} & {{ft_K_U30F_vs_CTRL_h100}} {{ft_K_U30F_vs_CTRL_h100_ci}} & {{ft_K_U30F_vs_CTRL_h368}} {{ft_K_U30F_vs_CTRL_h368_ci}} \\
U50F & control & {{ft_K_U50F_vs_CTRL_h1}} {{ft_K_U50F_vs_CTRL_h1_ci}} & {{ft_K_U50F_vs_CTRL_h24}} {{ft_K_U50F_vs_CTRL_h24_ci}} & {{ft_K_U50F_vs_CTRL_h100}} {{ft_K_U50F_vs_CTRL_h100_ci}} & {{ft_K_U50F_vs_CTRL_h368}} {{ft_K_U50F_vs_CTRL_h368_ci}} \\
\bottomrule
\end{tabular}}
\end{table}

\textbf{What was known before the registration.} Table~\ref{tab:decisions}'s $K$ for the 30\%-pruned copies without
fine-tuning (seeds 0--4) was known: {{K_u30_h1}} at one step, {{K_u30_h100}} at 100 steps and {{K_u30_h368}} at 368
steps. So were the earlier unit-removal study's registered results (Table~\ref{tab:app-perc}): fine-tuned width-192
copies kept most of the control's 368-step skill ({{perc_K_pool_w192_h368}}, pooled), width 128 was not settled, and
fine-tuning alone took the unreduced control to {{perc_ctrl_pool_h368}} of the original's skill at 368 steps. No
fine-tuned magnitude-pruned copy had been trained beyond a few test iterations, and none had been scored; $K$ of the
unretrained copies for seeds {{ft_seeds}} alone could have been computed from stored scores, but the registration
records that it was not.
